\clearpage
\section{辛对合}\label{sec:14}\label{section:14}
\subsection{等变膜}\label{sec:14a}
设 $(M,\allowbreak \omega_M)$ 为辛流形，配有殆复结构 $I_M$、二次复体积形式 $\eta_M^2$，以及保持 $(\omega_M,\allowbreak I_M,\allowbreak \eta_M^2)$ 的对合 $\iota_M$。\znote{199}{此处 $\iota_M$ 称为辛对合：它是二阶辛同胚，且保持 $I_M$，故在每点切空间上给出复线性对合，其不动点集 $M^{\mathrm{fix}}$ 因而是复子流形，法向切空间带复结构；当法向复维数为一（最简单的非平凡情形）时 $M^{\mathrm{fix}}$ 的实余维为二，这就是后文假设 14.3 的几何来源；保持 $\eta_M^2$ 则保证不动点集上仍留有可用的二次复体积形式，后文讨论分次时正要用到它。若法向复维数更高，下述整套论证需要更多技术条件，故假设 14.3 只处理余维二的情形。[来源：原书 §(14c) 假设 14.3，印刷页 203。]}希望将 $\iota_M$ 对 Lagrangian 子流形的自然作用与 Pin 结构的改变结合。这在纯形式层面造成一些困难，下面予以处理。
\begin{definition*}
\emph{预等变膜}由以下数据组成：其主体，即 Lagrangian 膜 $L_+^\#=(L_+,\allowbreak \alpha_+^\#,\allowbreak P_+^\#)$；其影子，即分次 Lagrangian 子流形 $(L_-=\iota_M(L_+),\allowbreak \alpha_-^\#=\alpha_+^\#\circ\iota_M)$ 上的 Pin 结构 $P_-^\#$；以及一对同构
\[
f_+:P_+^\#\longrightarrow\iota_M^*(P_-^\#\otimes\lambda^{\mathrm{top}}(TL_-)),\qquad
f_-:P_-^\#\longrightarrow\iota_M^*(P_+^\#\otimes\lambda^{\mathrm{top}}(TL_+)),
\]
使复合
\begin{equation}\label{eq:14.1}
\begin{aligned}
P_+^\#&\xrightarrow{f_+}\iota_M^*(P_-^\#\otimes\lambda^{\mathrm{top}}(TL_-))\\
&\xrightarrow{\iota_M^*(f_-)\otimes\Id}
P_+^\#\otimes\lambda^{\mathrm{top}}(TL_+)\otimes\iota_M^*\lambda^{\mathrm{top}}(TL_-)
\cong P_+^\#
\end{aligned}
\end{equation}
等于 Pin 结构的非平凡自同构 $-e\in\ker(\operatorname{Pin}_n\to O_n)$；最后一步使用显然的同构 $\iota_M^*(TL_-)\cong TL_+$。\znote{199}{$\operatorname{Pin}_n\to O_n$ 是二重覆叠，核同构于 $\Z/2$，其非平凡元素记作 $-e$（$e$ 为恒等元），它实现标架丛二重覆叠的换层。按二次型 $\pm\sum x_i^2$ 的选取，Pin 群有两种约定 $\operatorname{Pin}_\pm(n)$，二者相差一个符号；因此“平方等于 $e$ 还是 $-e$”必须逐处指明，定义中特意写 $-e$ 而不写“平凡”，正是为了固定这一约定。当 $n$ 较小时两种约定还可能给出同构的群，但一般情形下并不相同，因此对照其他文献时必须先核对它采取哪一种，否则符号会整体差一个。[来源：nLab "pin group"，https:/\allowbreak /\allowbreak ncatlab.org/\allowbreak nlab/\allowbreak show/\allowbreak pin+\allowbreak group（2026-11-15 访问）]}
\end{definition*}
为使措辞简洁，有时称之为“预等变膜 $L_+^\#$”，但总默认辅助数据 $P_-^\#$ 与 $f_-$ 已作某个具体选择；也可用 $f_+$，因为两个同构中任一个都唯一确定另一个。
\begin{note}\label{st:14.1}
每个 Lagrangian 膜 $L_+^\#$ 都可以显然地成为预等变膜：令 $P_-^\#=\iota_M^*(P_+^\#\otimes\lambda^{\mathrm{top}}(TL_+))$，取 $f_-$ 为恒等，并相应调整 $f_+$。
\end{note}
预等变膜上有自然对合 $\epsilon$，交换正、负两部分；简记为 $\epsilon(L_+^\#)=L_-^\#$。这里使用了另一顺序的复合 $(\iota_M^*(f_+)\otimes\Id)\circ f_-$ 也等于 $-e$ 这一事实。
\begin{definition*}
\emph{等变膜}是此对合的不动点。具体地，它由不变 Lagrangian 子流形 $L$、$\iota_M(L)=L$，一个 $\iota_M$ 不变分次 $\alpha^\#$，以及 Pin 结构 $P^\#$ 连同同构 $f:P^\#\to\iota_M^*(P^\#\otimes\lambda^{\mathrm{top}}(TL))$ 组成，使其“平方”
\SourcePage{207}
\begin{equation}\label{eq:14.2}
\begin{aligned}
P^\#&\xrightarrow{f}\iota_M^*(P^\#\otimes\lambda^{\mathrm{top}}(TL))\\
&\xrightarrow{\iota_M^*(f)\otimes\Id}P^\#\otimes\lambda^{\mathrm{top}}(TL)\otimes\iota_M^*\lambda^{\mathrm{top}}(TL)\cong P^\#
\end{aligned}
\end{equation}
等于 $-e$。
\end{definition*}
给定 $\iota_M$ 不变 Lagrangian 子流形 $L$，考察它容许等变膜结构的条件。假设 $H^1(L)$ 中的 Maslov 类消失，使 $L$ 容许分次 $\alpha^\#$。若 $L$ 连通，则差 $\alpha^\#\circ\iota_M-\alpha^\#$ 是 $\iota_M$ 反不变的常值函数，因此自动为零。另一方面，若 $L$ 有两个被 $\iota_M$ 交换的连通分支，可调整其中一个分支上的 $\alpha^\#$ 使该差消失。一般情形由这两种情形得到；结论是，得到不变分次没有额外障碍。

Pin 部分更有趣。最直接地，可以逐步处理：先为 $L$ 赋予 Pin 结构，再试图找到同构 $f:P^\#\to\iota_M^*(P^\#\otimes\lambda^{\mathrm{top}}(TL))$，最后确定它的平方。障碍分别位于 $H^2(L;\Z/2)$、$H^1(L;\Z/2)$ 和 $H^0(L;\Z/2)$；但后两项是次级障碍，依赖先前选择，并不取遍所有值。也可将这些方面结合成一个更典范的单一障碍，即等变上同调类
\begin{equation}\label{eq:14.3}
w_2^{\mathrm{eq}}(L)+w_1^{\mathrm{eq}}(L)u+u^2\in H_{\Z/2}^2(L;\Z/2).
\end{equation}
这里 $u$ 为 $H^*(B\Z/2;\Z/2)=\Z/2[u]$ 的生成元，$w_k^{\mathrm{eq}}(L)$ 为 $TL$ 的等变 Stiefel–Whitney 类，依赖其等变向量丛结构。若 $\iota_M|L$ 没有不动点，有一个稳定化技巧，可将问题归结到非等变 Pin 结构的标准问题，从而说明\eref{14.3}。具体地，令 $\pi:L\to\bar L$ 为到商 $\bar L=L/(\Z/2)$ 的覆盖，$\beta\to\bar L$ 为关联实线丛。\znote{199}{$\Z/2$ 在 $L$ 上自由作用，故商 $\bar L=L/(\Z/2)$ 是光滑流形、$\pi$ 是二重覆叠；$\beta$ 是符号表示所关联的实线丛，其拉回 $\pi^*\beta$ 平凡。于是 $TL$ 的等变结构与 $T\bar L\oplus\beta$ 的结构相当，再添三个 $\beta$ 使作用自由，就把等变 Pin 结构化为通常的 Pin 结构；$\R^3$ 的维数三正是引理 11.14 中 $\operatorname{rank}(F^{\mathrm{anti}})=3$ 情形的来源。反过来看，正因为 $\pi^*\beta$ 平凡、且 $\Z/2$ 通过符号表示作用在 $\beta$ 上，$L$ 上的 Pin 结构与 $\bar L$ 上的 Pin 结构加 $\beta$ 的扭曲才只是同一份数据的两种写法。[来源：原书 §(14a)，印刷页 199–200，并见引理 11.14。]}假设在 $T\bar L\oplus\beta\oplus\beta\oplus\beta$ 上有 Pin 结构 $\bar P^\#$，即
\begin{equation}\label{eq:14.4}
\begin{aligned}
w_2(T\bar L\oplus\beta\oplus\beta\oplus\beta)
&=w_2(T\bar L)+w_1(T\bar L)w_1(\beta)+w_1(\beta)^2\\
&\in H^2(\bar L;\Z/2)
\end{aligned}
\end{equation}
必须为零。取拉回 $\pi^*\bar P^\#$，它是 $TL\oplus\R^3$ 上的 Pin 结构，令 $P^\#$ 为其在 $TL$ 上的限制。该拉回有显然的自同构，覆盖 $D(\iota_M|L)\oplus(-\Id_{\R^3})$ 在标架丛上的作用。希望将其与覆盖 $\Id_{TL}\oplus(-\Id_{\R^3})$ 的辅助自同构 $g$ 复合，再将复合限制到 $P^\#$。由\eref{11.30}与\stref{引理}{11.14}，$g$ 的选择对应于丛 $\R^3\otimes\lambda^{\mathrm{top}}(TL\oplus\R^3)$ 的定向，也就是 $TL$ 本身的定向。将这一因子计入后，例如选择局部定向，再拼接所得 $P^\#$ 的局部自同构，就得到 $P^\#$ 的自然自同构 $f:P^\#\to\iota_M^*(P^\#\otimes\lambda^{\mathrm{top}}(TL))$。依定义，$\pi^*\bar P^\#$ 的原始自同构为对合；由
\SourcePage{208}
\stref{引理}{11.14}中 $\operatorname{rank}(F^{\mathrm{anti}})=3$ 的情形，有 $g^2=-e$；二者交换，因而\eref{14.2}具有所需性质。论证也可反向进行，说明等变膜结构的存在蕴含\eref{14.4}消失，而在 $H_{\Z/2}^*(L)\cong H^*(\bar L)$ 下，它恰对应于\eref{14.3}。

对一般的、非自由作用的情形，可以通过 Borel 构造将任意 $\Z/2$ 作用转为自由作用，立即说明\eref{14.3}消失是 $L$ 上等变膜结构存在的必要条件。\znote{200}{Borel 构造把任意 $\Z/2$ 作用换成自由作用：取 $L\times_{\Z/2}E\Z/2$（对角作用），它是同伦商，其普通上同调即等变上同调 $H^*_{\Z/2}(L)$；当 $\Z/2$ 自由作用时它同伦等价于商空间 $L/(\Z/2)$。因此“等变上同调类消失”可以在这一自由作用空间上按通常的示性类理论验证，这便是正文用它论证必要性的道理。同伦商的构造不依赖于 $E\Z/2$ 的具体模型，故所得障碍类只依赖作用本身，可作为不变量使用。[来源：nLab "Borel construction"，https:/\allowbreak /\allowbreak ncatlab.org/\allowbreak nlab/\allowbreak show/\allowbreak Borel+\allowbreak construction（2026-11-15 访问）]}也可用这种方法证明充分性，但需要进一步讨论 Borel 构造上的 Pin 结构，这里省略。

为补充前述讨论，更直接地考察不变 Lagrangian 子流形 $L$ 的不动点集 $L^{\mathrm{fix}}$ 上的情形。将 $TL|L^{\mathrm{fix}}=F^{\mathrm{inv}}\oplus F^{\mathrm{anti}}$ 分成 $D\iota_M$ 不变与反不变部分。假设 $L$ 上有 Pin 结构 $P^\#$ 及同构 $f:P^\#\to\iota_M^*(P^\#\otimes\lambda^{\mathrm{top}}(TL))$。由\eref{11.30}，$f$ 在 $L^{\mathrm{fix}}$ 上的限制由
\[
\operatorname{Pin}(TL\otimes\lambda^{\mathrm{top}}(TL))\times_{\Z/2}S(\lambda^{\mathrm{top}}(TL))
\]
的一个截面给出，它覆盖 $D(\iota_M|L)\otimes\Id\in O(TL\otimes\lambda^{\mathrm{top}}(TL))$。再用\stref{引理}{11.14}可知，这种截面对应于 $(F^{\mathrm{anti}}\otimes\lambda^{\mathrm{top}}(TL))\oplus\lambda^{\mathrm{top}}(TL)$ 的定向，或等价地，对应于下列线丛的平凡化：
\begin{equation}\label{eq:14.5}
\lambda^{\mathrm{top}}(F^{\mathrm{anti}})^{\otimes\operatorname{rank}(F^{\mathrm{anti}})}
\otimes\lambda^{\mathrm{top}}(F^{\mathrm{inv}})^{\otimes(\operatorname{rank}(F^{\mathrm{anti}})+1)}\longrightarrow L^{\mathrm{fix}}.
\end{equation}
此外，若 $\frac12\operatorname{rank}(F^{\mathrm{anti}})(\operatorname{rank}(F^{\mathrm{anti}})+1)$ 为偶数，$f$ 的平方为 $e$；若为奇数，则为 $-e$。这里的平方按\eref{14.2}理解，所以相比\eref{11.29}多出符号 $(-1)^{\operatorname{rank}(F^{\mathrm{anti}})}$；它来自 $\det(D(\iota_M|L))$，是配对 $\iota_M^*(\lambda^{\mathrm{top}}(TL))\otimes\lambda^{\mathrm{top}}(TL)\to\R$ 出现的结果。顺便指出，这一论证的大部分也可从抽象障碍理论\eref{14.3}出发进行，只需注意，令 $r=\operatorname{rank}(F^{\mathrm{anti}})$、$F=TL|L^{\mathrm{fix}}$，有
\[
\begin{aligned}
&(w_2^{\mathrm{eq}}(L)+w_1^{\mathrm{eq}}(L)u+u^2)|L^{\mathrm{fix}}\\
&\quad=w_2(F)+(r\,w_1(F^{\mathrm{anti}})+(r+1)w_1(F^{\mathrm{inv}}))u
+\left(\frac{r(r+1)}2+1\right)u^2\\
&\quad\in H^*(L^{\mathrm{fix}};\Z/2)[u].
\end{aligned}
\]
最终，只有一个特殊情形与应用有关，将其记录下来以备后用：
\begin{lemma}\label{st:14.2}
设 $L$ 为 $M$ 中 $\iota_M$ 不变的 Lagrangian 子流形。假设 $L$ 连通，其不动点集 $L^{\mathrm{fix}}$ 也连通、非空且余维为一。再假设 $L$ 上有同构于 $\iota_M^*(P^\#\otimes\lambda^{\mathrm{top}}(TL))$ 的 Pin 结构 $P^\#$。则同构 $f:P^\#\to\iota_M^*(P^\#\otimes\lambda^{\mathrm{top}}(TL))$ 的选择与 $L^{\mathrm{fix}}\subset L$ 的余定向双射对应。此外，任一这样的同构按\eref{14.2}取平方，均为 $-e$。
\end{lemma}
\SourcePage{209}
显然，$M$ 的 $\Z/2$ 等变分次辛自同构作用于预等变膜集合，并与对合 $\epsilon$ 严格交换，特别地保持等变膜。等变线丛 $\beta\to M$ 也可作用：给定这种线丛及预等变膜，将 $P_\pm^\#$ 换成 $P_\pm^\#\otimes\beta|L_\pm$，利用 $\beta$ 的等变结构给出的同构 $\beta|L_\pm\cong\iota_M^*(\beta|L_\mp)$，将 $f_\pm$ 变成新的 Pin 结构对之间的同构。例如，取 $\beta=\R\times M$，带非平凡 $\Z/2$ 作用，对应于 $\Z/2$ 的非平凡实特征标；与它作张量积，等价于将两个 $f_\pm$ 均乘以 $-e$。回到一般 $\beta$，将此作用施加于等变膜，并考察不动点集 $L^{\mathrm{fix}}\subset L$，其效果是使\eref{14.5}的平凡化改变一个符号，该符号为 $\Z/2$ 在 $\beta|L^{\mathrm{fix}}$ 上的作用符号。
